\documentclass[
	fontsize=12pt,
	paper=a4,
	numbers=noenddot,
	parskip=half,
	openany
]{scrbook}
\usepackage{tcolorbox}

\include{../../../for_latex/preamble}
% ##################################################
% Gemeinsames Setup fuer die Probeklausuren
% (Java-Listings, Hinweis-/Beispiel-Boxen, Struktogramm-Stile in TikZ)
% wird nach \include{../for_latex/preamble} per \input eingebunden
% ##################################################

% --- Java Code-Highlighting ---
\definecolor{codegray}{rgb}{0.5,0.5,0.5}
\definecolor{codepurple}{rgb}{0.58,0,0.82}
\definecolor{backcolour}{rgb}{0.95,0.95,0.92}
\definecolor{keywordcolor}{rgb}{0.8,0.47,0.13}

\lstdefinestyle{javastyle}{
	backgroundcolor=\color{backcolour},
	commentstyle=\color{codegray},
	keywordstyle=\color{keywordcolor}\bfseries,
	numberstyle=\tiny\color{codegray},
	stringstyle=\color{codepurple},
	basicstyle=\ttfamily\footnotesize,
	breakatwhitespace=false,
	breaklines=true,
	captionpos=b,
	keepspaces=true,
	numbers=left,
	numbersep=5pt,
	showspaces=false,
	showstringspaces=false,
	showtabs=false,
	tabsize=2,
	language=Java,
	frame=single,
	rulecolor=\color{black!30}
}

\lstset{style=javastyle}

% --- Hinweis-/Beispiel-/Formel-Umgebungen (schlicht, ohne farbige Box) ---
\newenvironment{hinweis}[1][Hinweis]%
	{\par\smallskip\noindent\textbf{#1.}\enspace\ignorespaces}%
	{\par\smallskip}

\newenvironment{beispiel}[1][Beispiel]%
	{\par\smallskip\noindent\textbf{#1:}\par\nobreak\vspace{2pt}}%
	{\par\smallskip}

\newenvironment{formeln}[1][Gegebene Formeln]%
	{\par\smallskip\noindent\textbf{#1:}\par\nobreak\vspace{2pt}}%
	{\par\smallskip}

% ##################################################
% Struktogramm (Nassi-Shneiderman) in reinem TikZ
% ##################################################
\newlength{\nswidth}\setlength{\nswidth}{11.5cm}   % Gesamtbreite
\newlength{\nsind}\setlength{\nsind}{0.55cm}        % Einrueckung je Schleifenebene
% vorberechnete Versatzbreiten fuer Verzweigungen (in Koordinaten nutzbar)
\newlength{\nshalf}\setlength{\nshalf}{\dimexpr\nswidth/2\relax}            % halbe Gesamtbreite
\newlength{\nshalfi}\setlength{\nshalfi}{\dimexpr(\nswidth-\nsind)/2\relax}  % halbe Breite, 1x eingerueckt
\newlength{\nshalfii}\setlength{\nshalfii}{\dimexpr(\nswidth-2\nsind)/2\relax} % halbe Breite, 2x eingerueckt

\tikzset{
	% Basis fuer Anweisungs-/Verzweigungsbloecke
	nb/.style={draw, line width=0.4pt, anchor=north west, inner sep=4pt,
		align=left, font=\footnotesize, fill=white, minimum height=0.8cm},
	% volle Breite
	nbfull/.style={nb, minimum width=\nswidth,
		text width=\dimexpr\nswidth-8pt\relax},
	% eine Ebene eingerueckt (Schleifenrumpf)
	nbi/.style={nb, minimum width=\dimexpr\nswidth-\nsind\relax,
		text width=\dimexpr\nswidth-\nsind-8pt\relax},
	% zwei Ebenen eingerueckt (Schleife in Schleife)
	nbii/.style={nb, minimum width=\dimexpr\nswidth-2\nsind\relax,
		text width=\dimexpr\nswidth-2\nsind-8pt\relax},
	% Verzweigung: halbe Breiten
	nbh/.style={nb, minimum width=\dimexpr\nswidth/2\relax,
		text width=\dimexpr\nswidth/2-8pt\relax},
	nbih/.style={nb, minimum width=\dimexpr(\nswidth-\nsind)/2\relax,
		text width=\dimexpr(\nswidth-\nsind)/2-8pt\relax},
	% Verzweigung halbe Breite, zwei Ebenen eingerueckt
	nbiih/.style={nb, minimum width=\nshalfii,
		text width=\dimexpr\nshalfii-8pt\relax},
	% Kopf einer Verzweigung (das Dreieck), inner sep 0
	nc/.style={draw, line width=0.4pt, anchor=north west, inner sep=0pt,
		minimum height=1.0cm, fill=white},
	ncfull/.style={nc, minimum width=\nswidth},
	nci/.style={nc, minimum width=\dimexpr\nswidth-\nsind\relax},
	ncii/.style={nc, minimum width=\dimexpr\nswidth-2\nsind\relax},
}

% Zeichnet das Dreieck + Beschriftungen in einen bereits gesetzten nc-Knoten.
% #1 Knotenname  #2 Bedingung  #3 ja-Label  #4 nein-Label
\newcommand{\nscond}[4]{%
	\draw[line width=0.4pt] (#1.north west) -- (#1.south);
	\draw[line width=0.4pt] (#1.north east) -- (#1.south);
	\node[font=\scriptsize, anchor=north, inner sep=1pt] at ($(#1.north)+(0,-1.5pt)$) {#2};
	\node[font=\scriptsize, anchor=south west, inner sep=1pt] at ($(#1.south west)+(2.5pt,1.5pt)$) {#3};
	\node[font=\scriptsize, anchor=south east, inner sep=1pt] at ($(#1.south east)+(-2.5pt,1.5pt)$) {#4};
}

% Zeichnet den L-Rahmen einer Schleife (linker + unterer Rand der Einrueckspalte).
% #1 Kopfknoten der Schleife  #2 letzter Rumpfknoten
\newcommand{\nsframe}[2]{%
	\draw[line width=0.4pt] (#1.south west) |- (#2.south west);
}


\title{Probeklausur 1}
\author{Luis Staudt}
\date{}

\begin{document}

% --- Titelblock ---
\begin{center}
	\rule{\textwidth}{0.8pt}\\[0.6em]
	{\LARGE\bfseries Probeklausur 1}\\[0.3em]
	{\large Grundlagen der Programmierung}\\[0.7em]
	\begin{tabular}{r l @{\hspace{1.4cm}} r l}
		\textbf{Studiengänge:}    & PEB \& MVB            & \textbf{Autor:}        & Luis Staudt \\
		\textbf{Semester:}        & Sommersemester 2026   & \textbf{Schwierigkeit:}& einfach \\
		\textbf{Bearbeitungszeit:}& 90 Minuten            & \textbf{Gesamtpunkte:} & 90 \\
	\end{tabular}\\[0.7em]
	\rule{\textwidth}{0.8pt}
\end{center}
\vspace{0.5em}

\noindent\textit{Bearbeiten Sie alle fünf Aufgaben. Achten Sie auf sauberen, kompilierfähigen Java-Quelltext. Hilfsmittel: OpenBook, keine elektronischen Geräte.}

% --- Seitenkopf und -fuß ---
\pagestyle{fancy}
\fancyhf{}
\lhead{\small Probeklausur 1, Grundlagen der Programmierung}
\rhead{\small SoSe 2026 (einfach)}
\cfoot{\thepage}
\renewcommand{\headrulewidth}{0.4pt}

% =============================================
\clearpage
\section*{Aufgabe 1: Struktogramm in Java umsetzen \hfill (24 Punkte)}

Gegeben ist das folgende Struktogramm.
Es liest eine Anzahl von Messwerten ein, bestimmt deren größten Wert und Mittelwert und gibt anschließend zu jedem Wert aus, ob er unter oder über dem Mittelwert liegt.

Setzen Sie das Struktogramm \textbf{vollständig in ein lauffähiges Java-Programm} um.
Verwenden Sie zum Einlesen \texttt{Tastatureingabe.leseInt()} bzw.\\\ \texttt{Tastatureingabe.leseDouble()}.

\begin{center}
\begin{tikzpicture}
	\node[nbfull] (a1) at (0,0) {lies die Anzahl $n$ ein};
	\node[nbfull] (a2) at (a1.south west) {lege ein Feld \texttt{werte} mit $n$ Plätzen an};

	% Schleife 1
	\node[nbfull] (h1) at (a2.south west) {\textbf{für} $i = 0,\,1,\,\ldots,\,n-1$};
	\node[nbi]    (b1) at ($(h1.south west)+(\nsind,0)$) {lies \texttt{werte[i]} über die Tastatur ein};
	\nsframe{h1}{b1}

	\node[nbfull] (a3) at ($(b1.south west)+(-\nsind,0)$)
		{\texttt{summe = 0;}\qquad \texttt{max = werte[0];}};

	% Schleife 2
	\node[nbfull] (h2) at (a3.south west) {\textbf{für} $i = 0,\,1,\,\ldots,\,n-1$};
	\node[nci]    (c1) at ($(h2.south west)+(\nsind,0)$) {};
	\nscond{c1}{\texttt{werte[i] > max}}{ja}{nein}
	\node[nbih]   (c1l) at (c1.south west) {\texttt{max = werte[i];}};
	\node[nbih]   (c1r) at ($(c1.south west)+(\nshalfi,0)$) {\textit{(nichts)}};
	\node[nbi]    (b2) at (c1l.south west) {\texttt{summe = summe + werte[i];}};
	\nsframe{h2}{b2}

	\node[nbfull] (a4) at ($(b2.south west)+(-\nsind,0)$) {\texttt{mittel = summe / n;}};

	% Schleife 3
	\node[nbfull] (h3) at (a4.south west) {\textbf{für} $i = 0,\,1,\,\ldots,\,n-1$};
	\node[nci]    (c2) at ($(h3.south west)+(\nsind,0)$) {};
	\nscond{c2}{\texttt{werte[i] < mittel}}{ja}{nein}
	\node[nbih]   (c2l) at (c2.south west) {drucke \texttt{werte[i]}, \glqq unter\grqq};
	\node[nbih]   (c2r) at ($(c2.south west)+(\nshalfi,0)$) {drucke \texttt{werte[i]}, \glqq über\grqq};
	\nsframe{h3}{c2l}
\end{tikzpicture}
\end{center}

% =============================================
\clearpage
\section*{Aufgabe 2: Bewegungsformel auswählen \hfill (14 Punkte)}

Ein Körper bewegt sich geradlinig.
Je nach Bewegungsart beschreibt eine andere Formel die zurückgelegte Strecke~$s$.

\begin{formeln}
\[
	\text{(1) gleichförmig:}\quad s = v \cdot t
	\qquad
	\text{(2) aus Ruhe beschleunigt:}\quad s = \tfrac{1}{2}\,a\,t^2
\]
\[
	\text{(3) mit Anfangsgeschwindigkeit:}\quad s = v_0 \cdot t + \tfrac{1}{2}\,a\,t^2
\]
\end{formeln}

Schreiben Sie ein Programm, das den Benutzer zunächst die Bewegungsart per Zahl (\texttt{1}, \texttt{2} oder \texttt{3}) wählen lässt und anschließend die benötigten Größen abfragt.
Berechnen und drucken Sie die Strecke~$s$ mit der \textbf{richtigen} Formel. Fangen Sie eine \textbf{negative Zeit} (\texttt{t < 0}) als ungültige Eingabe ab.
Die Auswahl ist mit \texttt{if/else} (kein \texttt{switch} nötig) zu realisieren.

\begin{beispiel}[Beispieldialog]
\begin{lstlisting}[language={}]
Bewegungsart (1=gleichfoermig, 2=beschleunigt, 3=mit v0): 2
Beschleunigung a [m/s^2]: 2.5
Zeit t [s]: 4
Zurueckgelegte Strecke s = 20.0 m
\end{lstlisting}
\end{beispiel}

% =============================================
\clearpage
\section*{Aufgabe 3: Punkt im Kreis \hfill (20 Punkte)}

Gegeben ist die aus Vorlesung und Übung bekannte Klasse \texttt{Punkt} (Attribute \texttt{xKoordinate}, \texttt{yKoordinate} als \texttt{double}).

Implementieren Sie die Methode \texttt{liegtInnerhalb}, die zurückgibt, ob die aufrufende Punktinstanz innerhalb eines Kreises mit Radius~$r$ um einen übergebenen Mittelpunkt liegt.
Übergabeparameter sind der Radius sowie der Mittelpunkt (ein \texttt{Punkt}).

\begin{center}
\begin{tabular}{c l}
	\toprule
	\textbf{Rückgabe} & \textbf{Bedeutung} \\
	\midrule
	\texttt{1}  & Punkt liegt innerhalb des Kreises oder auf dem Umfang \\
	\texttt{0}  & Punkt liegt außerhalb des Kreises \\
	\texttt{-1} & ungültiger Kreis (Radius~$\le 0$) \\
	\bottomrule
\end{tabular}
\end{center}

\begin{hinweis}[Wiederverwendung]
	Wer die aus der Übung bekannte Methode\\ \texttt{berechneAbstandZuPunkt(Punkt p)} sinnvoll wiederverwendet, erhält einen \textbf{Zusatzpunkt}.
\end{hinweis}

\begin{lstlisting}
Punkt p1 = new Punkt();
p1.xKoordinate = 1.0;  p1.yKoordinate = 1.5;
Punkt mp = new Punkt();
mp.xKoordinate = 5.6;  mp.yKoordinate = -4.3;

int i = p1.liegtInnerhalb(3.0, mp);
if (i == 1)       System.out.println("p1 liegt innerhalb");
else if (i == 0)  System.out.println("p1 liegt ausserhalb");
else              System.out.println("ungueltiger Radius");
\end{lstlisting}

% =============================================
\clearpage
\section*{Aufgabe 4: Fachwerk, Stäbe klassifizieren \hfill (16 Punkte)}

Gegeben ist die aus Vorlesung und Übung bekannte Datei \texttt{fachwerk.jar}. Nach der Berechnung soll für ein gewähltes Fachwerk gezählt werden, wie viele \textbf{Zugstäbe}, \textbf{Druckstäbe} und \textbf{Nullstäbe} es enthält.

\begin{center}
\begin{tikzpicture}[>=Stealth, scale=1.0,
	knoten/.style={circle, fill=black, inner sep=1.6pt},
	stab/.style={thick, gray!70!black}]
	\coordinate (A) at (0,0);
	\coordinate (B) at (2.5,0);
	\coordinate (C) at (5,0);
	\coordinate (D) at (1.25,1.9);
	\coordinate (E) at (3.75,1.9);
	\draw[stab] (A)--(B); \draw[stab] (B)--(C);
	\draw[stab] (A)--(D); \draw[stab] (D)--(B);
	\draw[stab] (D)--(E); \draw[stab] (E)--(B); \draw[stab] (E)--(C);
	\foreach \p in {A,B,C,D,E} \node[knoten] at (\p) {};
	\draw[->, thick] (B)--++(0,-0.9);
	\node[below, font=\footnotesize] at ($(B)+(0,-0.9)$) {Last};
\end{tikzpicture}\\[-0.3em]
{\footnotesize\itshape Schematische Darstellung eines Fachwerks.}
\end{center}

\begin{hinweis}[Nullstab über Epsilon-Umgebung]
	Die Berechnung liefert für Nullstäbe sehr kleine Werte ($\ne 0$), z.\,B. \texttt{3.55E-15}. Erkennen Sie einen Nullstab über eine \textbf{Epsilon-Umgebung}:
	\[ |F| < \varepsilon \;\Rightarrow\; \text{Nullstab}, \qquad
	   F > \varepsilon \;\Rightarrow\; \text{Zugstab}, \qquad
	   F < -\varepsilon \;\Rightarrow\; \text{Druckstab.} \]
	Benötigte Methoden: \texttt{int gibAnzahlStaebe()}, \texttt{Fachwerkstab gibStab(int i)}, \texttt{double gibStabkraft()}.
\end{hinweis}

Ergänzen Sie den folgenden Quelltext, sodass die drei Anzahlen ausgegeben werden:

\begin{lstlisting}
public class Fachwerkaufgabe {
    public static void main(String[] args) {
        System.out.println("Welches Fachwerk wollen Sie rechnen (0-4)?");
        int wahl = Tastatureingabe.leseInt();
        Fachwerk fw = new Fachwerk(wahl);
        fw.berechneGroessen();
        double epsilon = 1e-6;
        // hier erweitern: zaehlen und ausgeben
    }
}
\end{lstlisting}

% =============================================
\clearpage
\section*{Aufgabe 5: Wertetabelle in eine Datei schreiben \hfill (16 Punkte)}

Ein Körper fällt aus der Ruhe; für die zurückgelegte Strecke gilt
\[ s(t) = \tfrac{1}{2}\,a\,t^2 \qquad \text{mit } a = 9{,}81~\text{m/s}^2. \]

Schreiben Sie ein Programm, das für die Zeitpunkte $t = 0~\text{s}$ bis $t = 15~\text{s}$ in Schritten von $1~\text{s}$ jeweils die Strecke berechnet und Zeile für Zeile in die Datei \texttt{fall.txt} schreibt.

\begin{hinweis}[Hinweis]
	Verwenden Sie \texttt{java.io.FileWriter}: Datei mit \lstinline|new FileWriter("fall.txt")| öffnen, je Zeile mit \texttt{write(...)} schreiben (Zeilenumbruch \lstinline|"\n"| nicht vergessen) und am Ende \texttt{close()} aufrufen.
	Nutzen Sie \lstinline|"\t"| für die Spaltenabstände.
\end{hinweis}

\begin{beispiel}[Erwarteter Dateiinhalt (Auszug)]
\begin{lstlisting}[language={}]
t = 0 s    s = 0.0 m
t = 1 s    s = 4.905 m
t = 2 s    s = 19.62 m
...
t = 15 s   s = 1103.625 m
\end{lstlisting}
\end{beispiel}

\end{document}
